\documentclass[11pt,a4paper]{article}
\usepackage[utf8]{inputenc}
\usepackage[T1]{fontenc}
\usepackage[spanish,es-minimal,es-noshorthands]{babel}
\usepackage{lmodern}
\usepackage[sfdefault,lining]{FiraSans}
\usepackage[a4paper,top=2.3cm,bottom=2.2cm,left=1.8cm,right=1.8cm,headheight=1.3cm,headsep=0.4cm,footskip=1cm]{geometry}
\usepackage{amsmath,amssymb}
\usepackage{icomma}
\usepackage{graphicx}
\usepackage[table]{xcolor}
\usepackage{tikz}
\usepackage{fontawesome5}
\usepackage{ccicons}
\usepackage{fancyhdr}
\usepackage{lastpage}
\usepackage{tasks}
\usepackage{needspace}
\usepackage[most]{tcolorbox}
\usepackage{array}
\usepackage[hidelinks]{hyperref}
\AtBeginDocument{\hypersetup{pdftitle={Soluciones de la ficha de la unidad 2: Polinomios y fracciones algebraicas. 4º ESO},pdfsubject={Matemáticas B 4º ESO},pdfkeywords={polinomios, valor numérico, identidades notables, binomio de Newton, Ruffini, teorema del resto, factorización, m.c.d., m.c.m., fracciones algebraicas}}}

\definecolor{navy}{HTML}{1B2748}
\definecolor{accent}{HTML}{E8603A}
\definecolor{teal}{HTML}{1F9E90}
\definecolor{soft}{HTML}{EEF1F8}
\definecolor{muted}{HTML}{5D6685}

\setlength{\parindent}{0pt}
\setlength{\parskip}{3pt}
\settasks{label=\alph*), label-format=\bfseries\color{navy}, label-width=1.3em, item-indent=1.8em, column-sep=0.6em, after-item-skip=5pt, before-skip=3pt, after-skip=4pt}

% ---------- capçalera i peu ----------
\newcommand{\llicencia}{\href{https://creativecommons.org/licenses/by-nc-sa/4.0/deed.es}{\ccbyncsa}\ \ Licencia Creative Commons Atribución-NoComercial-CompartirIgual 4.0 Internacional (CC BY-NC-SA 4.0)}
\pagestyle{fancy}
\fancyhf{}
\fancyhead[L]{\raisebox{-0.35\height}{\includegraphics[height=0.95cm]{logo-gva.png}}}
\fancyhead[C]{\small\color{muted}Soluciones de la ficha · Unidad 2: Polinomios y fracciones algebraicas · 4\textsuperscript{o} ESO B}
\fancyhead[R]{\raisebox{-0.35\height}{\includegraphics[height=0.95cm]{logo-vermellar.png}}}
\fancyfoot[L]{\scriptsize\color{muted}\llicencia}
\fancyfoot[R]{\small\color{muted}\thepage}
\renewcommand{\headrulewidth}{0.4pt}
\renewcommand{\headrule}{\color{muted!40}\hrule height \headrulewidth}
\fancypagestyle{primera}{\fancyhf{}\renewcommand{\headrulewidth}{0pt}
  \fancyfoot[L]{\scriptsize\color{muted}\llicencia}
  \fancyfoot[R]{\small\color{muted}\thepage}}

% ---------- nivells i exercicis ----------
\newcommand{\estrelles}[1]{{\color{accent}\ifcase#1\or\faStar\faStar[regular]\faStar[regular]\or\faStar\faStar\faStar[regular]\or\faStar\faStar\faStar\fi}}
\newcounter{exer}
\newcommand{\ex}[1]{\par\needspace{4\baselineskip}\medskip\refstepcounter{exer}\noindent{\color{navy}\bfseries\theexer.}\ \ {\small\estrelles{#1}}\ \ }

\newcommand{\seccio}[2]{\par\needspace{16\baselineskip}\bigskip
  \begin{tcolorbox}[colback=navy,colframe=navy,coltext=white,arc=2mm,boxrule=0pt,left=3mm,right=3mm,top=1.2mm,bottom=1.2mm]
  {\large\bfseries #1\quad #2}
  \end{tcolorbox}\vspace{-2pt}}

\newcommand{\exf}[1]{\par\medskip\refstepcounter{exer}\noindent{\color{navy}\bfseries\theexer.}\ \ {\small\estrelles{#1}}\ \ }
\newcolumntype{C}[1]{>{\centering\arraybackslash}p{#1}}
% quadre de recordatori (només a la fitxa d'exercicis)
\newcommand{\recorda}[1]{\par\smallskip\begin{tcolorbox}[colback=accent!7,colframe=accent!55,arc=2mm,boxrule=0.5pt,left=2.5mm,right=2.5mm,top=1mm,bottom=1mm]\small{\color{accent}\faLightbulb[regular]}\ \textbf{Recuerda:}\ #1\end{tcolorbox}}
\newcommand{\res}[1]{\boxed{#1}}\newcommand{\rest}[1]{\fbox{#1}}

\def\fmneed{11}% que el títol d'una secció no quede sol al final del full
% =====================================================================
% Capçalera FentMates de les fitxes: un color per a cada curs.
% Els colors són els de la pàgina inicial (LEVELS a assets/js/materials.js);
% el logo ve de _eines/logo/pdf/ (python3 _eines/logo/genera.py).
%
% Ús, al final del preàmbul de la unitat:
%   % =====================================================================
% Capçalera FentMates de les fitxes: un color per a cada curs.
% Els colors són els de la pàgina inicial (LEVELS a assets/js/materials.js);
% el logo ve de _eines/logo/pdf/ (python3 _eines/logo/genera.py).
%
% Ús, al final del preàmbul de la unitat:
%   % =====================================================================
% Capçalera FentMates de les fitxes: un color per a cada curs.
% Els colors són els de la pàgina inicial (LEVELS a assets/js/materials.js);
% el logo ve de _eines/logo/pdf/ (python3 _eines/logo/genera.py).
%
% Ús, al final del preàmbul de la unitat:
%   \input{../comu/fentmates.tex}
%   \fmcurs{1eso}
% i, al començament del document:
%   \capcalera{Matemàtiques}{1r ESO}{Fitxa d'exercicis · Unitat 2}{Divisibilitat}
% Opcional abans de \input: \def\fmneed{14} (espai mínim abans de cada secció).
% =====================================================================
\makeatletter
\def\fm@logodir{../../logo/pdf}
\providecommand{\fmneed}{10}
\def\fm@cursnom{}

% clau del curs, color i color del text damunt del color
\newcommand{\fm@curs}[3]{\@namedef{fm@color@#1}{#2}\@namedef{fm@text@#1}{#3}}
\fm@curs{1eso}{E6B84A}{navy}
\fm@curs{2eso}{3B82F6}{white}
\fm@curs{3eso}{FF7A50}{white}
\fm@curs{4eso-a}{A78BFA}{white}
\fm@curs{4eso-b}{6D28D9}{white}
\fm@curs{1batx-mi}{9B1C1C}{white}
\fm@curs{1batx-ccss}{EF6B6B}{white}
\fm@curs{2batx-mii}{1E6B47}{white}
\fm@curs{2batx-ccss}{2BB3A3}{white}

\newcommand{\fmcurs}[1]{%
  \@ifundefined{fm@color@#1}{\PackageError{fentmates}{Curs desconegut: #1}{}}{}%
  \edef\fm@tmp{\noexpand\definecolor{curs}{HTML}{\@nameuse{fm@color@#1}}%
               \noexpand\colorlet{curstext}{\@nameuse{fm@text@#1}}}\fm@tmp
  \colorlet{cursfosc}{curs!72!black}%
  \def\fm@logo{\fm@logodir/vertical-#1.pdf}%
  \def\fm@marca{\fm@logodir/marca-#1.pdf}}

% ---------- capçalera de la primera pàgina ----------
% #1 matèria (petit), #2 curs, #3 tipus de document i unitat, #4 títol
\newcommand{\capcalera}[4]{%
  \gdef\fm@cursnom{#2}%
  \par\noindent
  \begin{tcolorbox}[colback=curs,colframe=curs,boxrule=0pt,arc=2.2mm,width=3.5cm,height=2.15cm,
      valign=center,left=3mm,right=2mm,top=1mm,bottom=1mm,nobeforeafter,box align=center]
    \color{curstext}{\footnotesize #1}\par\vspace{1pt}{\fontsize{19}{21}\selectfont\bfseries #2\par}%
  \end{tcolorbox}\hspace{5mm}%
  \begin{minipage}[c]{\dimexpr\linewidth-3.5cm-5mm-2.9cm\relax}
    {\color{cursfosc}\small\bfseries\MakeUppercase{#3}}\par\vspace{2pt}
    {\color{navy}\fontsize{20}{23}\selectfont\bfseries #4\par}
  \end{minipage}\hfill
  \begin{minipage}[c]{2.6cm}\raggedleft\includegraphics[width=2.6cm]{\fm@logo}\end{minipage}\par
  \vspace{2pt}}

% ---------- barres de secció del color del curs ----------
\renewcommand{\seccio}[2]{\par\needspace{\fmneed\baselineskip}\bigskip
  \begin{tcolorbox}[colback=curs,colframe=curs,coltext=curstext,arc=2mm,boxrule=0pt,left=3mm,right=3mm,top=1.2mm,bottom=1.2mm]
  {\large\bfseries #1\quad #2}
  \end{tcolorbox}\vspace{-2pt}}

% ---------- estrelles de dificultat del color del curs (a Batxillerat no se'n posen) ----------
\@ifundefined{estrelles}{}{%
  \renewcommand{\estrelles}[1]{{\color{curs}\ifcase#1\or\faStar\faStar[regular]\faStar[regular]\or\faStar\faStar\faStar[regular]\or\faStar\faStar\faStar\fi}}}

% ---------- capçalera i peu de les pàgines ----------
% primera pàgina: logos institucionals al peu
\fancypagestyle{primera}{\fancyhf{}\renewcommand{\headrulewidth}{0pt}
  \fancyfoot[L]{\raisebox{-0.32\height}{\includegraphics[height=0.78cm]{logo-gva.png}}\hspace{5mm}%
    \raisebox{-0.32\height}{\includegraphics[height=0.78cm]{logo-vermellar.png}}}
  \fancyfoot[R]{\scriptsize\color{muted}\href{https://creativecommons.org/licenses/by-nc-sa/4.0/deed.ca}{\ccbyncsa}\ \ CC BY-NC-SA 4.0 · fentmates.org\qquad\small\thepage}}
% la resta: el símbol a l'esquerra i el curs a la dreta, amb el títol al centre com abans
\fancyhead[L]{\raisebox{-0.3\height}{\includegraphics[height=0.62cm]{\fm@marca}}}
\fancyhead[R]{\small\bfseries\color{cursfosc}\fm@cursnom}
\renewcommand{\headrule}{\color{curs!45}\hrule height \headrulewidth}
\makeatother

%   \fmcurs{1eso}
% i, al començament del document:
%   \capcalera{Matemàtiques}{1r ESO}{Fitxa d'exercicis · Unitat 2}{Divisibilitat}
% Opcional abans de \input: \def\fmneed{14} (espai mínim abans de cada secció).
% =====================================================================
\makeatletter
\def\fm@logodir{../../logo/pdf}
\providecommand{\fmneed}{10}
\def\fm@cursnom{}

% clau del curs, color i color del text damunt del color
\newcommand{\fm@curs}[3]{\@namedef{fm@color@#1}{#2}\@namedef{fm@text@#1}{#3}}
\fm@curs{1eso}{E6B84A}{navy}
\fm@curs{2eso}{3B82F6}{white}
\fm@curs{3eso}{FF7A50}{white}
\fm@curs{4eso-a}{A78BFA}{white}
\fm@curs{4eso-b}{6D28D9}{white}
\fm@curs{1batx-mi}{9B1C1C}{white}
\fm@curs{1batx-ccss}{EF6B6B}{white}
\fm@curs{2batx-mii}{1E6B47}{white}
\fm@curs{2batx-ccss}{2BB3A3}{white}

\newcommand{\fmcurs}[1]{%
  \@ifundefined{fm@color@#1}{\PackageError{fentmates}{Curs desconegut: #1}{}}{}%
  \edef\fm@tmp{\noexpand\definecolor{curs}{HTML}{\@nameuse{fm@color@#1}}%
               \noexpand\colorlet{curstext}{\@nameuse{fm@text@#1}}}\fm@tmp
  \colorlet{cursfosc}{curs!72!black}%
  \def\fm@logo{\fm@logodir/vertical-#1.pdf}%
  \def\fm@marca{\fm@logodir/marca-#1.pdf}}

% ---------- capçalera de la primera pàgina ----------
% #1 matèria (petit), #2 curs, #3 tipus de document i unitat, #4 títol
\newcommand{\capcalera}[4]{%
  \gdef\fm@cursnom{#2}%
  \par\noindent
  \begin{tcolorbox}[colback=curs,colframe=curs,boxrule=0pt,arc=2.2mm,width=3.5cm,height=2.15cm,
      valign=center,left=3mm,right=2mm,top=1mm,bottom=1mm,nobeforeafter,box align=center]
    \color{curstext}{\footnotesize #1}\par\vspace{1pt}{\fontsize{19}{21}\selectfont\bfseries #2\par}%
  \end{tcolorbox}\hspace{5mm}%
  \begin{minipage}[c]{\dimexpr\linewidth-3.5cm-5mm-2.9cm\relax}
    {\color{cursfosc}\small\bfseries\MakeUppercase{#3}}\par\vspace{2pt}
    {\color{navy}\fontsize{20}{23}\selectfont\bfseries #4\par}
  \end{minipage}\hfill
  \begin{minipage}[c]{2.6cm}\raggedleft\includegraphics[width=2.6cm]{\fm@logo}\end{minipage}\par
  \vspace{2pt}}

% ---------- barres de secció del color del curs ----------
\renewcommand{\seccio}[2]{\par\needspace{\fmneed\baselineskip}\bigskip
  \begin{tcolorbox}[colback=curs,colframe=curs,coltext=curstext,arc=2mm,boxrule=0pt,left=3mm,right=3mm,top=1.2mm,bottom=1.2mm]
  {\large\bfseries #1\quad #2}
  \end{tcolorbox}\vspace{-2pt}}

% ---------- estrelles de dificultat del color del curs (a Batxillerat no se'n posen) ----------
\@ifundefined{estrelles}{}{%
  \renewcommand{\estrelles}[1]{{\color{curs}\ifcase#1\or\faStar\faStar[regular]\faStar[regular]\or\faStar\faStar\faStar[regular]\or\faStar\faStar\faStar\fi}}}

% ---------- capçalera i peu de les pàgines ----------
% primera pàgina: logos institucionals al peu
\fancypagestyle{primera}{\fancyhf{}\renewcommand{\headrulewidth}{0pt}
  \fancyfoot[L]{\raisebox{-0.32\height}{\includegraphics[height=0.78cm]{logo-gva.png}}\hspace{5mm}%
    \raisebox{-0.32\height}{\includegraphics[height=0.78cm]{logo-vermellar.png}}}
  \fancyfoot[R]{\scriptsize\color{muted}\href{https://creativecommons.org/licenses/by-nc-sa/4.0/deed.ca}{\ccbyncsa}\ \ CC BY-NC-SA 4.0 · fentmates.org\qquad\small\thepage}}
% la resta: el símbol a l'esquerra i el curs a la dreta, amb el títol al centre com abans
\fancyhead[L]{\raisebox{-0.3\height}{\includegraphics[height=0.62cm]{\fm@marca}}}
\fancyhead[R]{\small\bfseries\color{cursfosc}\fm@cursnom}
\renewcommand{\headrule}{\color{curs!45}\hrule height \headrulewidth}
\makeatother

%   \fmcurs{1eso}
% i, al començament del document:
%   \capcalera{Matemàtiques}{1r ESO}{Fitxa d'exercicis · Unitat 2}{Divisibilitat}
% Opcional abans de \input: \def\fmneed{14} (espai mínim abans de cada secció).
% =====================================================================
\makeatletter
\def\fm@logodir{../../logo/pdf}
\providecommand{\fmneed}{10}
\def\fm@cursnom{}

% clau del curs, color i color del text damunt del color
\newcommand{\fm@curs}[3]{\@namedef{fm@color@#1}{#2}\@namedef{fm@text@#1}{#3}}
\fm@curs{1eso}{E6B84A}{navy}
\fm@curs{2eso}{3B82F6}{white}
\fm@curs{3eso}{FF7A50}{white}
\fm@curs{4eso-a}{A78BFA}{white}
\fm@curs{4eso-b}{6D28D9}{white}
\fm@curs{1batx-mi}{9B1C1C}{white}
\fm@curs{1batx-ccss}{EF6B6B}{white}
\fm@curs{2batx-mii}{1E6B47}{white}
\fm@curs{2batx-ccss}{2BB3A3}{white}

\newcommand{\fmcurs}[1]{%
  \@ifundefined{fm@color@#1}{\PackageError{fentmates}{Curs desconegut: #1}{}}{}%
  \edef\fm@tmp{\noexpand\definecolor{curs}{HTML}{\@nameuse{fm@color@#1}}%
               \noexpand\colorlet{curstext}{\@nameuse{fm@text@#1}}}\fm@tmp
  \colorlet{cursfosc}{curs!72!black}%
  \def\fm@logo{\fm@logodir/vertical-#1.pdf}%
  \def\fm@marca{\fm@logodir/marca-#1.pdf}}

% ---------- capçalera de la primera pàgina ----------
% #1 matèria (petit), #2 curs, #3 tipus de document i unitat, #4 títol
\newcommand{\capcalera}[4]{%
  \gdef\fm@cursnom{#2}%
  \par\noindent
  \begin{tcolorbox}[colback=curs,colframe=curs,boxrule=0pt,arc=2.2mm,width=3.5cm,height=2.15cm,
      valign=center,left=3mm,right=2mm,top=1mm,bottom=1mm,nobeforeafter,box align=center]
    \color{curstext}{\footnotesize #1}\par\vspace{1pt}{\fontsize{19}{21}\selectfont\bfseries #2\par}%
  \end{tcolorbox}\hspace{5mm}%
  \begin{minipage}[c]{\dimexpr\linewidth-3.5cm-5mm-2.9cm\relax}
    {\color{cursfosc}\small\bfseries\MakeUppercase{#3}}\par\vspace{2pt}
    {\color{navy}\fontsize{20}{23}\selectfont\bfseries #4\par}
  \end{minipage}\hfill
  \begin{minipage}[c]{2.6cm}\raggedleft\includegraphics[width=2.6cm]{\fm@logo}\end{minipage}\par
  \vspace{2pt}}

% ---------- barres de secció del color del curs ----------
\renewcommand{\seccio}[2]{\par\needspace{\fmneed\baselineskip}\bigskip
  \begin{tcolorbox}[colback=curs,colframe=curs,coltext=curstext,arc=2mm,boxrule=0pt,left=3mm,right=3mm,top=1.2mm,bottom=1.2mm]
  {\large\bfseries #1\quad #2}
  \end{tcolorbox}\vspace{-2pt}}

% ---------- estrelles de dificultat del color del curs (a Batxillerat no se'n posen) ----------
\@ifundefined{estrelles}{}{%
  \renewcommand{\estrelles}[1]{{\color{curs}\ifcase#1\or\faStar\faStar[regular]\faStar[regular]\or\faStar\faStar\faStar[regular]\or\faStar\faStar\faStar\fi}}}

% ---------- capçalera i peu de les pàgines ----------
% primera pàgina: logos institucionals al peu
\fancypagestyle{primera}{\fancyhf{}\renewcommand{\headrulewidth}{0pt}
  \fancyfoot[L]{\raisebox{-0.32\height}{\includegraphics[height=0.78cm]{logo-gva.png}}\hspace{5mm}%
    \raisebox{-0.32\height}{\includegraphics[height=0.78cm]{logo-vermellar.png}}}
  \fancyfoot[R]{\scriptsize\color{muted}\href{https://creativecommons.org/licenses/by-nc-sa/4.0/deed.ca}{\ccbyncsa}\ \ CC BY-NC-SA 4.0 · fentmates.org\qquad\small\thepage}}
% la resta: el símbol a l'esquerra i el curs a la dreta, amb el títol al centre com abans
\fancyhead[L]{\raisebox{-0.3\height}{\includegraphics[height=0.62cm]{\fm@marca}}}
\fancyhead[R]{\small\bfseries\color{cursfosc}\fm@cursnom}
\renewcommand{\headrule}{\color{curs!45}\hrule height \headrulewidth}
\makeatother

\fmcurs{4eso-b}

\begin{document}
\thispagestyle{primera}

\capcalera{Matemáticas B}{4\textsuperscript{o} ESO}{Soluciones de la ficha · Unidad 2}{\raggedright Polinomios y fracciones algebraicas}
\vspace{8pt}

\begin{tcolorbox}[colback=soft,colframe=soft,arc=2mm,left=3mm,right=3mm,top=1.5mm,bottom=1.5mm]
\small Intenta resolver cada ejercicio antes de mirar su solución. Puedes comprobar las divisiones, Ruffini, las raíces, las factorizaciones y las fracciones con las herramientas interactivas de los apuntes de la unidad.
\end{tcolorbox}

% =====================================================
\seccio{1}{Polinomios y valor numérico}

\exf{1} \begin{center}\renewcommand{\arraystretch}{1.45}\small\begin{tabular}{|l|C{1cm}|C{1.9cm}|C{1.5cm}|C{2cm}|C{1.7cm}|}\hline\rowcolor{soft}\textbf{Polinomio} & \textbf{Grado} & \textbf{Térm. indep.} & \textbf{N.\textsuperscript{o} términos} & \textbf{¿Completo?} & \textbf{Coef. princ.}\\\hline$P(x)=5x^4-3x^3+2x+10$ & $\res{4}$ & $\res{10}$ & $\res{4}$ & \rest{Incompleto} & $\res{5}$\\\hline$Q(x)=-5x^5-\frac32x^4-x^3+9x^2-6$ & $\res{5}$ & $\res{-6}$ & $\res{5}$ & \rest{Incompleto} & $\res{-5}$\\\hline$R(x)=4x^2+x+1$ & $\res{2}$ & $\res{1}$ & $\res{3}$ & \rest{Completo} & $\res{4}$\\\hline$S(x)=5x^2+2x^5+9x$ & $\res{5}$ & $\res{0}$ & $\res{3}$ & \rest{Incompleto} & $\res{2}$\\\hline$T(x)=3x-2$ & $\res{1}$ & $\res{-2}$ & $\res{2}$ & \rest{Completo} & $\res{3}$\\\hline$U(x)=x^3-4x^2+x+1$ & $\res{3}$ & $\res{1}$ & $\res{4}$ & \rest{Completo} & $\res{1}$\\\hline\end{tabular}\end{center}

\ex{1} $P(2)=-48-16-4-1=\res{-69}$; \quad $P(-2)=-48+16-4-1=\res{-37}$.

\ex{1} \begin{tasks}(2)
  \task $\res{150},\ \ \res{-10},\ \ \res{-3}$
  \task $\res{-124},\ \ \res{8},\ \ \res{5}$
  \task $\res{8},\ \ \res{8},\ \ \res{8}$
  \task $\res{26},\ \ \res{-2},\ \ \res{-1}$
\end{tasks}

\ex{2} $P(2)=16-8+10+k=18+k=6\Rightarrow\res{k=-12}$.

\ex{2} $P(-3)=-81+108-3a+10=37-3a=10\Rightarrow\res{a=9}$.


% =====================================================
\seccio{2}{Operaciones e identidades notables}

\exf{2} \begin{tasks}(1)
  \task $(-5x^{3}+4x^{2}-8)(4x^2-x)=\res{-20x^{5}+21x^{4}-4x^{3}-32x^{2}+8x}$
  \task $(x^{3}-2x+1)(2x^3-x+3)=\res{2x^{6}-5x^{4}+5x^{3}+2x^{2}-7x+3}$
\end{tasks}

\ex{1} \begin{tasks}(3)
  \task $\res{x^{2}+6x+9}$
  \task $\res{25x^{4}+70x^{2}+49}$
  \task $\res{4x^{2}-12x+9}$
  \task $\res{4x^{2}-25}$
  \task $\res{2x^{2}-20x+50}$
  \task $\res{\frac{4}{9}x^{2}-2x+\frac{9}{4}}$
\end{tasks}

\ex{2} \begin{tasks}(1)
  \task $\res{6x^{2}-7x+13}$
  \task $\res{-30x-77}$
  \task $\res{-\frac{4}{3}x^{3}-\frac{5}{3}x^{2}-8x+\frac{2}{3}}$
\end{tasks}


% =====================================================
\seccio{3}{Binomio de Newton}

\exf{2} \begin{tasks}(1)
  \task $\res{x^{3}+6x^{2}+12x+8}$
  \task $\res{x^{4}-4x^{3}+6x^{2}-4x+1}$
  \task $\res{8x^{3}+36x^{2}+54x+27}$
  \task $\res{32x^{5}+80x^{4}+80x^{3}+40x^{2}+10x+1}$
  \task $\res{8x^{3}-36x^{2}+54x-27}$
  \task $\res{x^{6}-12x^{5}+60x^{4}-160x^{3}+240x^{2}-192x+64}$
\end{tasks}

\ex{3} \begin{tasks}(1)
  \task $\res{x^{6} + 3 x^{5} + 3 x^{4} + x^{3}}$
  \task $\res{32 x^{20} - 80 x^{16} y + 80 x^{12} y^{2} - 40 x^{8} y^{3} + 10 x^{4} y^{4} - y^{5}}$
  \task $\res{256 x^{12} + 512 x^{9} + 384 x^{6} + 128 x^{3} + 16}$
  \task $\res{x^{30} - 15 x^{24} y^{2} + 90 x^{18} y^{4} - 270 x^{12} y^{6} + 405 x^{6} y^{8} - 243 y^{10}}$
  \task $\res{125 x^{3} + 225 x^{2} y^{2} + 135 x y^{4} + 27 y^{6}}$
  \task $\res{16 x^{12} - 64 x^{10} + 96 x^{8} - 64 x^{6} + 16 x^{4}}$
  \task $\res{64 x^{18} + 576 x^{16} + 2160 x^{14} + 4320 x^{12} + 4860 x^{10} + 2916 x^{8} + 729 x^{6}}$
  \task $\res{x^{4} - 8 x^{3} y + 24 x^{2} y^{2} - 32 x y^{3} + 16 y^{4}}$
\end{tasks}(En f) y g) conviene sacar primero el signo: $(-2x^3-3x)^6=(2x^3+3x)^6$.)


% =====================================================
\seccio{4}{División de polinomios y regla de Ruffini}

\exf{2} \begin{tasks}(1)
  \task $c(x)=\res{3x^{2}+8x+6},\quad r(x)=\res{9x+6}$
  \task $c(x)=\res{x^{3}-4x^{2}+3x+2},\quad r(x)=\res{-2x-7}$
  \task $c(x)=\res{4x^{3}-x+7},\quad r(x)=\res{7x+1}$
  \task $c(x)=\res{x^{3}-5},\quad r(x)=\res{0}$
  \task $c(x)=\res{18x^{4}+33x^{2}+2x+66},\quad r(x)=\res{5x+132}$
\end{tasks}

\ex{1} \begin{tasks}(1)
  \task $c(x)=\res{x^{2}-3x+2},\quad r(x)=\res{1}$
  \task $c(x)=\res{x^{3}+3x^{2}+x-5},\quad r(x)=\res{3}$
  \task $c(x)=\res{x^{6}-4x^{5}-x^{4}+7x^{3}+32x^{2}+179x+858},\quad r(x)=\res{4253}$
\end{tasks}

\ex{1} \begin{tasks}(3)
  \task $\res{-1}$
  \task $\res{41}$
  \task $\res{2}$
\end{tasks}


% =====================================================
\seccio{5}{Valor numérico y teorema del resto}

\exf{2} Debe ser $P(-1)=0$: $-1-m-5-2=0\Rightarrow\res{m=-8}$.

\ex{2} $P(2)=32+8k-14+6=24+8k=-8\Rightarrow\res{k=-4}$.

\ex{1} Sustituyendo: $48-4+4-5=\res{43}$. Con Ruffini, el resto de $P(x):(x-2)$ también es 43 (teorema del resto).

\ex{2} \begin{tasks}(2)
  \task $P(1)=2\neq0$: $\res{\text{no}}$
  \task $P(-1)=-1+1=0$: $\res{\text{sí}}$
\end{tasks}

\ex{2} El resto de dividir entre $x-3$ es $P(3)=0$: $\res{\text{sí lo es}}$.

\ex{2} El resto de dividir entre $x+4$ es $P(-4)=1\neq0$: $\res{\text{no lo es}}$.


% =====================================================
\seccio{6}{Raíces de un polinomio}

\exf{2} \begin{tasks}(2)
  \task $\res{x=1,\ x=2,\ x=3}$
  \task $\res{x=1,\ x=-1,\ x=-2}$
  \task $\res{x=\pm1,\ x=\pm3}$
  \task $\res{x=2,\ x=-2,\ x=3}$
\end{tasks}

\ex{2} Divisors de 12: $\pm1,\ \pm2,\ \pm3,\ \pm4,\ \pm6,\ \pm12$. Solo $P(2)=0$: la única raíz es $\res{x=2}$ (el cociente $x^2+3x+6$ no tiene raíces reales).

\ex{2} $(x-1)(x+2)(x-3)=\res{x^{3}-2x^{2}-5x+6}$; \ $(x-4)^2=\res{x^{2}-8x+16}$ (o cualquier múltiplo).


% =====================================================
\seccio{7}{Factorización}

\exf{1} \begin{tasks}(2)
  \task $\res{5y(xy+2)}$
  \task $\res{4x^2(3x-1)}$
  \task $\res{6(b^2-4)=6(b+2)(b-2)}$
  \task $\res{2ab^2c^3(2c-14abc^2+21a^4b^4)}$
  \task $\res{12(2z^2-z+1)}$
  \task $\res{5w^4t^4(t^2+5w^2)}$
\end{tasks}

\ex{2} \begin{tasks}(2)
  \task $\res{\left(x - 3\right) \left(x - 1\right) \left(x + 1\right) \left(x + 2\right)}$
  \task $\res{\left(x - 1\right)^{3} \left(x + 2\right)^{2}}$
  \task $\res{\left(x - 1\right)^{2} \left(x + 3\right)^{2}}$
  \task $\res{x \left(2 x^{4} - 6 x^{3} - 19 x^{2} + 24 x + 32\right)}$
  \task $\res{\left(x + 2\right) \left(x^{2} + x + 3\right)}$
  \task $\res{\left(x - 3\right) \left(x - 1\right) \left(x + 1\right) \left(2 x - 1\right)}$
\end{tasks}

\ex{2} \begin{tasks}(2)
  \task $\res{\left(x + 1\right) \left(x^{2} + x + 1\right)}$
  \task $\res{\left(x + 2\right) \left(x^{2} + 2 x + 2\right)}$
  \task $\res{\left(x - 1\right) \left(x^{2} + x + 2\right)}$
  \task $\res{\left(x - 2\right) \left(x^{2} + 2 x + 3\right)}$
  \task $\res{\left(x - 5\right) \left(x - 2\right) \left(x + 2\right) \left(x + 5\right)}$
  \task $\res{\left(x - 5\right) \left(x + 5\right) \left(x^{2} + 1\right)}$
  \task $\res{\left(x - 1\right) \left(x + 1\right) \left(x^{2} - 3 x + 5\right)}$
  \task $\res{\left(x - 2\right) \left(x + 2\right) \left(x^{2} + 2 x + 5\right)}$
\end{tasks}


% =====================================================
\seccio{8}{Divisibilidad. M.c.d. y m.c.m.}

\exf{1} \begin{tasks}(2)
  \task $1-1=0$: $\res{\text{sí}}$
  \task $4+4=8\neq0$: $\res{\text{no}}$
  \task $1-1=0$: $\res{\text{sí}}$
  \task $9-9=0$: $\res{\text{sí}}$
\end{tasks}

\ex{2} \begin{tasks}(1)
  \task m.c.d. $=\res{x+2}$, \ m.c.m. $=\res{(x+2)^2(x-2)}$
  \task m.c.d. $=\res{x-1}$, \ m.c.m. $=\res{x(x+1)(x-1)}$
  \task m.c.d. $=\res{x-1}$, \ m.c.m. $=\res{x(x+1)(x-1)^2}$
  \task m.c.d. $=\res{(x-1)(x-2)}$, \ m.c.m. $=\res{(x-1)(x+1)(x-2)}$
\end{tasks}


% =====================================================
\seccio{9}{Simplificación de fracciones algebraicas}

\exf{2} \begin{tasks}(3)
  \task $\res{\frac{\left(x + 1\right) \left(x + 2\right)}{x - 2}}$
  \task $\res{\frac{x^{2} + x + 1}{x + 1}}$
  \task $\res{\frac{x - 3}{x + 3}}$
  \task $\res{\frac{1}{\left(x - 3\right) \left(x + 3\right)}}$
  \task $\res{\frac{x^{2}}{x - 1}}$
  \task $\res{\frac{x^{2} + 1}{\left(x - 1\right) \left(x^{2} + x + 1\right)}}$
  \task $\res{\frac{x - 7}{x + 7}}$
  \task $\res{x}$
\end{tasks}

\ex{2} \begin{tasks}(3)
  \task $\res{\frac{x \left(2 x + 1\right)}{x + 3}}$
  \task $\res{x - 2}$
  \task $\res{\frac{\left(2 x - 5\right) \left(3 x + 4\right)}{2 \left(x - 3\right) \left(x + 3\right)}}$
  \task $\res{\frac{x^{2} + 4}{x + 4}}$
  \task $\res{\frac{x + 1}{x - 1}}$
\end{tasks}


% =====================================================
\seccio{10}{Multiplicación y división de fracciones algebraicas}

\exf{2} \begin{tasks}(3)
  \task $\res{\frac{x}{\left(x + 1\right) \left(x + 6\right)}}$
  \task $\res{\frac{x - 2}{x \left(x + 3\right)}}$
  \task $\res{\frac{\left(x + 4\right) \left(x + 5\right)}{x - 5}}$
  \task $\res{\frac{1}{2}}$
  \task $\res{\frac{x^{3} \left(x + 2\right)}{\left(x - 2\right) \left(x + 3\right)}}$
  \task $\res{\frac{1}{\left(x + 1\right)^{2}}}$
  \task $\res{\frac{x^{2}}{8}}$
  \task $\res{1}$
  \task $\res{x - 2}$
\end{tasks}


% =====================================================
\seccio{11}{Suma y resta de fracciones algebraicas}

\exf{2} \begin{tasks}(2)
  \task $\res{\frac{x^{2} + 6 x - 2}{x - 2}}$
  \task $\res{- \frac{2 x^{2} - 5 x - 9}{3 x + 2}}$
  \task $\res{- \frac{1}{\left(x - 3\right) \left(x - 1\right)}}$
  \task $\res{\frac{2 x}{\left(x - 5\right) \left(x + 5\right)}}$
  \task $\res{- \frac{2}{x + 2}}$
  \task $\res{- \frac{x^{3} - 2 x^{2} - 3 x - 3}{\left(x - 1\right) \left(x + 1\right) \left(x^{2} + x + 1\right)}}$
  \task $\res{\frac{2 \left(x^{2} + 1\right)}{\left(x - 1\right) \left(x + 1\right)}}$
  \task $\res{\frac{x - 1}{3 \left(x - 4\right)}}$
  \task $\res{\frac{5 x^{3} + 16 x^{2} + 4 x - 18}{x \left(x - 3\right) \left(x + 3\right)}}$
  \task $\res{\frac{x + 2}{x \left(x + 1\right)}}$
\end{tasks}

\ex{3} \begin{tasks}(2)
  \task $\res{\frac{3}{x + 1}}$
  \task $\res{\frac{2 x^{3} - 5 x^{2} + 9 x + 12}{\left(x - 3\right) \left(x + 3\right)}}$
  \task $\res{- \frac{2 x^{2} + 2 x - 3}{3 x \left(x + 1\right)}}$
  \task $\res{- \frac{x - 2}{x}}$
\end{tasks}


% =====================================================
\seccio{12}{Operaciones combinadas}

\exf{3} \begin{tasks}(2)
  \task $\res{\frac{a \left(a + 4 b\right)}{\left(a - b\right) \left(a + b\right)}}$
  \task $\res{\frac{x + 1}{x}}$
  \task $\res{a b \left(a - b\right)}$
  \task $\res{- \frac{x + y}{x - y}}$
  \task $\res{- \frac{x \left(x - 5\right)}{\left(x - 2\right) \left(x + 2\right)}}$
\end{tasks}


% =====================================================
\seccio{13}{Problemas}

\exf{2} Altura: $(x^2+5x+6):(x+3)=\res{x+2\ \text{cm}}$. Perímetro: $2(x+3)+2(x+2)=\res{4x+10\ \text{cm}}$.

\ex{2} $\dfrac{x(x+2)(x+4)}{x(x+2)}=\res{x+4\ \text{cm}}$.


\end{document}
